\documentclass[11pt]{article}
\usepackage[a4paper,margin=1in]{geometry}
\usepackage{amsmath,amssymb,mathtools,bm}
\usepackage{enumitem,booktabs,array}
\usepackage{tcolorbox}
\usepackage{hyperref}
\usepackage[protrusion=true,expansion=false]{microtype}
\hypersetup{colorlinks=true,linkcolor=blue,urlcolor=blue,citecolor=blue}
\setlist[itemize]{topsep=2pt,itemsep=2pt}
\setlist[enumerate]{topsep=2pt,itemsep=2pt}
\newtcolorbox{keybox}{colback=blue!3!white,colframe=blue!40!black,boxrule=0.4pt,arc=1mm}
\newtcolorbox{alertbox}{colback=red!3!white,colframe=red!40!black,boxrule=0.4pt,arc=1mm}
\newtcolorbox{answerbox}{colback=green!3!white,colframe=green!40!black,boxrule=0.4pt,arc=1mm}
\newtcolorbox{drillbox}{colback=yellow!5!white,colframe=orange!60!black,boxrule=0.4pt,arc=1mm}
\newcommand{\EV}{\mathbb{E}}
\newcommand{\Var}{\mathrm{Var}}
\newcommand{\Cov}{\mathrm{Cov}}
\newcommand{\blank}[1]{\underline{\hspace{#1}}}

\title{E03-01: Faleye, Hoitash \& Hoitash (2011, JFE) \\
\large ``The costs of intense board monitoring'' --- Active Recall Workbook}
\author{Empirical Corporate Finance II --- Exam Prep}
\date{\today}
\begin{document}\maketitle

\begin{keybox}
\textbf{Paper in one sentence.} Boards on which a majority of independent directors serve on all three monitoring committees (audit, compensation, nominating) --- labelled ``monitoring-intensive'' --- produce higher financial-statement quality but are associated with lower R\&D, fewer acquisitions, and lower firm value, especially in innovation-dependent firms, showing that monitoring has real costs as well as benefits.

\textbf{Placement.} Canonical paper on the \emph{trade-off} view of board monitoring vs.\ advising. Contrast with Sarbanes-Oxley's monitoring-heavy prescription. Paired with Masulis-Mobbs (E03-04) on director effort allocation.
\end{keybox}

\section*{Round 1: Complete Walk-Through}

\subsection*{1.1 Measurement of monitoring intensity}
A board is ``monitoring-intensive'' (MIB) if a majority of independent directors sit on all three principal monitoring committees: audit, compensation, and nominating/governance. This operationalizes Adams-Ferreira (2007) monitoring-advising trade-off.

\subsection*{1.2 Sample and spec}
\begin{itemize}
\item US listed firms, 1998--2006.
\item Board-committee data from proxy filings + IRRC.
\item Main spec: firm FE, year FE, firm controls; $Y_{it}=\alpha_i+\delta_t+\beta\cdot \text{MIB}_{it}+X_{it}'\gamma+\varepsilon_{it}$.
\item Outcomes: earnings-management proxies, R\&D, acquisitions, Tobin's Q.
\end{itemize}

\subsection*{1.3 Headline findings: benefits}
\begin{itemize}
\item MIB firms have \emph{lower} discretionary accruals (better reporting).
\item Lower probability of restatements.
\item CEO pay is more performance-sensitive.
\item Consistent with MIB boards exercising monitoring function.
\end{itemize}

\subsection*{1.4 Headline findings: costs}
\begin{itemize}
\item R\&D intensity lower at MIB firms; innovation output (patents) lower.
\item Fewer acquisitions; when they occur, lower announcement returns.
\item Tobin's Q lower, especially in innovation-heavy industries.
\item Intuition: monitoring-intense boards crowd out advising and long-term risk-taking.
\end{itemize}

\subsection*{1.5 Heterogeneity}
\begin{itemize}
\item Costs are concentrated in R\&D-intensive and young firms.
\item Benefits are widespread but modest in size.
\item Net effect varies by firm type: MIB helps transparency-needy firms, hurts innovation-needy firms.
\end{itemize}

\subsection*{1.6 Policy implications}
\begin{itemize}
\item One-size-fits-all monitoring rules (e.g.\ SOX) can harm innovation-focused firms.
\item Board design should match firm informational needs.
\item Motivates a board-type-heterogeneity view of governance.
\end{itemize}

\subsection*{1.7 Caveats}
\begin{alertbox}
\begin{itemize}
\item MIB is endogenous: firms choose committee structure.
\item Some identification via within-firm variation from post-SOX adjustments; but no clean IV.
\item Patent-count measures of innovation are noisy.
\end{itemize}
\end{alertbox}

\section*{Round 2: Level 1 Fill-in}
\begin{enumerate}
\item Monitoring-intensive board: majority of independent directors on all \blank{1cm} monitoring committees.
\item Benefit: lower discretionary \blank{2cm} and fewer restatements.
\item Cost: lower \blank{1cm}\&D, fewer \blank{2cm}, lower Tobin's Q.
\item Costs concentrated in \blank{2cm}-heavy / young firms.
\item Policy lesson: \blank{2cm}-fits-all monitoring rules can harm innovation.
\end{enumerate}

\section*{Round 3: Level 2 Prompts}
\begin{enumerate}
\item Define monitoring-intensive boards and explain the monitoring-advising trade-off.
\item Describe the spec and the main empirical findings (benefits and costs).
\item Discuss the heterogeneity result and its policy implications.
\item Relate FHH to Adams-Ferreira (2007) and to SOX.
\item What endogeneity issues arise, and how does the paper address them?
\end{enumerate}

\section*{Round 4: Minimal Scaffolding}
From a blank page: (i) MIB definition, (ii) benefits, (iii) costs, (iv) heterogeneity, (v) policy.

\section*{Round 5: Blank Page}
Essay: monitoring vs.\ advising --- the costs of intense board monitoring.

\vspace{6pt}
\begin{drillbox}
\textbf{Speed Drill.} (1) MIB definition. (2) Accounting benefit. (3) Innovation cost. (4) Heterogeneity group. (5) Policy lesson.
\end{drillbox}

\section*{Quick Reference \& Answer Key}
\begin{answerbox}
\begin{itemize}
\item \textbf{Definition.} Majority indep. on all 3 monitoring committees.
\item \textbf{Benefit.} Better reporting.
\item \textbf{Cost.} Lower R\&D, fewer M\&A, lower Q.
\item \textbf{Heterogeneity.} Costs bite in innovation-heavy firms.
\end{itemize}
\end{answerbox}

\begin{keybox}
\textbf{30-second exam pitch.} Faleye, Hoitash \& Hoitash (2011) operationalize the board monitoring--advising trade-off of Adams-Ferreira (2007): a ``monitoring-intensive board'' (MIB) has a majority of independent directors on all three monitoring committees (audit, compensation, nominating). Using US panel data 1998--2006 with firm and year fixed effects, MIB firms produce \emph{better} financial-statement quality (lower discretionary accruals, fewer restatements) and more-performance-sensitive CEO pay --- the monitoring benefit --- but \emph{worse} real outcomes: lower R\&D, fewer acquisitions, lower Tobin's Q, concentrated in R\&D-intensive and young firms. Intense monitoring crowds out advising and risk-taking needed for innovation. The paper cautions against one-size-fits-all monitoring mandates (SOX) and supports firm-type-specific board design. FHH is the empirical benchmark for the cost-of-monitoring thesis and pairs naturally with Masulis-Mobbs (2014) on director effort allocation.
\end{keybox}

\end{document}
